\section{兼顾高亏格的 GW 理论导论（Patrick Lei）}%
\label{sec:Intro to GW theory with a view to higher genus (Patrick Lei)}


\subsection{GW 理论简介}%
\label{sub:Introduction to GW theory}

关于虚基本类的构造，参见~\cite{BehrendFantechi1997}；关于公理与量子上同调，参见~\cite{KM1994}。

这些讲座的目标是以一种便于进行高亏格计算的方式讨论格罗莫夫–威滕理论。格罗莫夫–威滕理论研究模空间 \(\Mbar_{g,n}(X, \beta)\) 上的积分，其中 \(X\) 是光滑射影代数簇，\(\beta \in H_2(X, \Z)\) 是有效曲线类。有如下图表
\begin{equation*}
\begin{tikzcd}
    \mc{C} \ar{r}{f} \ar[bend right=20]{d}[swap]{\pi} & X \\
    \Mbar_{g,n}(X, \beta) \ar[bend right=20]{u}[swap]{\sigma_i} \ar{ur}[swap]{\on{ev}_i}
\end{tikzcd}
\end{equation*}
其中 \(\mc{C}\) 是万有曲线，\(\sigma_i\) 是给出标记点的 \(\pi\) 的截面，\(f\) 是万有稳定映射，\(\ev_i = f \circ \sigma_i\) 是求值态射。尽管 \(\Mbar_{g,n}(X, \beta)\) 一般是奇异的、非等维的且可约的，但存在一个特殊的类
\[ [\Mbar_{g,n}(X, \beta)]^{\vir} \in A_{\ab(c_1(X), \beta) + (\dim X - 3)(1-g) + n}(\Mbar_{g,n}(X, \beta), \Q), \]
它称为\textit{虚基本类}。

\begin{exer}
    利用形变理论中的标准论证，验证 \(\Mbar_{g,n}(X, \beta)\) 的“期望维数”就是上面给出的数。
\end{exer}

在 \(X\) 光滑射影的假设下，\(\Mbar_{g,n}(X, \beta)\) 是真（proper）的，故可定义\textit{GW 不变量}
\[ \ab<\tau_{a_1}(\alpha_1), \cdots, \tau_{a_n}(\alpha_n)>_{g,n,\beta} \coloneqq \int_{[\Mbar_{g,n}(X, \beta)]^{\vir}} \prod_{i=1}^n \ev_i^* \alpha_i \cdot \psi_i^{a_i} \]
其中 \(\alpha_i \in H^*(X)\)。这里 \(\psi_i\) 由 \(\psi_i = c_1(\sigma_i^*(\omega_{\pi}))\) 定义。

遗憾的是，计算单个格罗莫夫–威滕不变量非常困难，因此我们转而考虑更为困难的问题：研究 GW 不变量的生成级数，例如当 \(X\) 为卡拉比–丘三维簇时的
\[ F_g = \sum_{\beta \geq 0} q^{\beta} \ab<>_{g,0,\beta} \]
这里 \(q\) 是记录曲线类的形式变量，我们将环 \(\Q\ps{q^{\beta} \mid \beta \in H_2(X, \Z)_{\mr{eff}}}\) 记作 \(\Lambda\)。我们还要考虑一些其他结构。首先可以考虑\textit{量子上同调环}，其乘积 \(a \star_{\mbf{t}} b\) 由下式给出
\[
    (a \star_{\mbf{t}} b, c) = \sum_{\substack{\beta \geq 0 \\ n \geq 0}} \ab<a,b,c, \mbf{t}, \ldots, \mbf{t}>_{0,3+n,\beta} \frac{q^{\beta}}{n!}.
\]
这里 \(a,b,c, \mbf{t} \in H^*(X)\)。我们还要考虑\textit{量子联络}。设 \(\ab\{ e_i \}\) 是 \(H^*(X)\) 的一组基。则量子联络是平凡丛上的联络，其纤维为 \(H^*(X)\)，底空间为 \(H^*(X) \times \P^1_z\)，由下式给出
\[ \nabla_{t^i} = \partial_{t^i} + \frac{1}{z} e_i \star_{\mbf{t}}. \]
有时考虑 \(z\) 方向上的联络也很有用，它由下式给出
\[ \nabla_{z \odv{}{z}} \coloneqq z \odv{}{z} - \frac{1}{z} E_X \star_{\mbf{t}} + \mu_X, \]
其中 \(\mu_X(\alpha) = \frac{\deg(\alpha) - \dim X}{2} \alpha\) 是分次算子，而
\[ E_X = c_1(X) + \sum_i \ab(1 - \frac{\deg(e_i)}{2})t^i e_i \]
是\textit{欧拉向量场}。由于我们主要关注卡拉比–丘簇，故不再考虑它。

\subsection{亏格零的枚举镜像对称}%
\label{sub:Enumerative mirror symmetry in genus zero}

数学家发展并研究格罗莫夫–威滕理论的动力来自 20 世纪 90 年代初：当时 Candelas、de la Ossa、Green 与 Parkes~\cite{CandelasEtAl1991} 通过把问题转化为计算另一个称为\textit{镜像五次三维簇}的卡拉比–丘三维簇上的周期积分（等价地，研究其霍奇结构的变分），预言了一般五次三维簇上低次数有理曲线的数目。在亏格零时，所有不变量均可由\textit{J-函数} \(J_{\mbf{t}}(q,z) \in H^*(X, \Lambda)[z]\) 恢复，其定义为
\[ \ab( J_{\mbf{t}}(q, z), \alpha ) \coloneqq (z + \mbf{t}, \alpha) + \sum_{n, \beta} \frac{q^{\beta}}{n!} \ab< \frac{\alpha}{z-\psi}, \mbf{t}, \ldots, \mbf{t}>_{0,1+n, \beta} \]
其中 \(\mbf{t} \in H^*(X, \Lambda)[z]\)，\(\alpha \in H^*(X)\)，而 \((-,-)\) 是庞加莱配对。因此，计算所有带后代（即 \(\psi\)-类）的亏格零格罗莫夫–威滕不变量的问题，归结为计算 J-函数。

\begin{rmk}
    我们将 \(\frac{1}{z-\psi}\) 展开为 \(\sum_{k=0}^{\infty} \psi^k z^{-k-1}\)。  
\end{rmk}

我们应当指出，J-函数可以通过量子微分方程的基本解得到，该基本解由下式给出
\[ ( S_{\mbf{t}}(z) \phi , \alpha) = (\phi, \alpha) + \sum_{n, \beta} \frac{q^{\beta}}{n!} \ab<\frac{\alpha}{z-\psi}, \phi, \mbf{t}, \ldots, \mbf{t}>_{0, 2+n, \beta}. \]
它具有性质 \(S_{\mbf{t}}^*(-z) S_{\mbf{t}}(z) = \on{Id}\)，这通常称为\textit{辛性质}。由此我们可由 \(z S_{\mbf{t}}^*(z)\cdot 1 = J_{\mbf{t}}(q, z)\) 恢复 J-函数。

在物理学中，镜像对称是称为 A 模型与 B 模型的两种理论之间的对偶，因此现在我们需要引入 B 模型的对象。为节省时间（并简化记号），我们假设 \(X\) 是环面簇且 \(-K_X\) 数值有效（nef）。设 \(D_1, \ldots, D_r\) 为 \(X\) 中全部的环面不变除子。定义\textit{\(I\)-函数}如下
\[
    I_{\mbf{t}}(q, z) \coloneqq z e^{\mbf{t}/z} \sum_{\beta \geq 0} q^{\beta} \frac{\prod_{i=1}^r \prod_{m = -\infty}^{0} (D_i + mz)}{\prod_{i=1}^r \prod_{m=-\infty}^{\beta \cdot D_i} (D_i + mz)}.
\]
注意，此处我们将出现的任意 \((D_i + mz)^{-1}\) 展开为 \(\frac{1}{mz} \ab(1 - \frac{D_i}{mz} + \cdots)\)，即展开为 \(\frac{D_i}{mz}\) 的几何级数。此外，I-函数被一个称为\textit{Picard–Fuchs 方程}的微分方程所零化，它是某种超几何型方程。

\begin{exm}
    若 \(\mbf{t} = 0\) 且 \(X = \P^1\)，则得到
    \begin{align*}
        I(q,z) &= z\sum_{d \geq 0} q^d \frac{1}{\prod_{m=1}^d (H + mz)^2} \\
        &= z\sum_{d \geq 0} \ab(\frac{q}{z^2})^d \frac{1}{(d!^2)} + O(H),
    \end{align*}
    于是实际上我们看到首项为第一类贝塞尔函数。
\end{exm}

若 \(X\) 是环面簇 \(Y\) 中由线丛 \(L_1, \ldots, L_k\) 切出的完全交，且这些线丛都满足 \(H^1(f^* L_i) = 0\)（对任意到 \(Y\) 的亏格零稳定映射成立，这称为\textit{凸}的），则量子 Lefschetz 定理~\cite{CoatesGivental2007} 允许我们把 \([\Mbar_{g,n}(X, \beta)]^{\vir}\) 替换为 \(e\ab(\bigoplus \pi_* f^* L_i) \cap [\Mbar_{g,n}(Y, \beta)]^{\vir}\)，并在 \(Y\) 上进行所有计算。此时 I-函数定义为
\[
    I_{\mbf{t}}(q,z) \coloneqq z e^{\mbf{t}/z} \sum_{\beta \geq 0} q^{\beta} \frac{\prod_{i=1}^k \prod_{m=-\infty}^{L_i \cdot \beta} (L_i + mz)}{\prod_{i=1}^k \prod_{m=-\infty}^{0} (L_i + mz)} \frac{\prod_{i=1}^r \prod_{m = -\infty}^{0} (D_i + mz)}{\prod_{i=1}^r \prod_{m=-\infty}^{\beta \cdot D_i} (D_i + mz)},
\]
其中我们滥用记号，将 \(L_i\) 与其第一陈类混同。

\begin{exm}
    对于由 \(\mc{O}_{\P^4}(5)\) 的一个截面切出的五次三维簇，\(\mbf{t} = 0\) 处的 I-函数为
    \[ I(q,z) = z \sum_{d \geq 0} q^d \frac{\prod_{m=1}^{5d}(5H+mz)}{\prod_{m=1}^d (H+mz)^5}. \]
\end{exm}

亏格零镜像定理的内容是这两个对象之间的对应。更精确地说，在 \(-K_X\) 数值有效的情形下，它们通过变量变换相联系。

\begin{thm}[{\cite{Givental1998,LLY1997}}]
    将 \(\mbf{t} = 0\) 的 I-函数写成 \(I_0(q) \cdot z + I_1(q) + O(z^{-1})\)。则有
    \[
        J_{\tau(\mbf{t})}(q,z) = \frac{I_{\mbf{t}}(q,z)}{I_0(q)},
    \]
    其中 \(\tau(\mbf{t}) = \frac{\mbf{t} + I_1(q)}{I_0(q)}\)。
\end{thm}

\begin{exm}
    当 \(X = \P^1\) 时，注意到 \(I_0(q) = 1\) 且 \(I_1(q) = 0\)，故实际上我们有等式
    \[ J_{\mbf{t}}(q,z) = I_{\mbf{t}}(q,z). \]
\end{exm}

\begin{exm}
    当 \(X\) 为五次三维簇时，为方便起见，我们将 I-函数的各分量改写为 \(I_0(q) \cdot z + I_1(q)H + O(z^{-1})\)。我们有
    \[ I_0(q) = \sum_{d \geq 0} \frac{(5d)!}{(d!)^5} q^d, \qquad I_1(q) = \sum_{d \geq 0} \frac{(5d)!}{(d!)^5} \ab(5\sum_{m=d+1}^{5d} \frac{1}{m}) q^d, \]
    故此情形下镜像映射非平凡。取 \(\mbf{t} = 0\)，我们看到 \(\tau = \frac{I_1(q)H}{I_0(q)}\)，于是得到
    \[ J_{\tau}(q,z) = \frac{I(q,z)}{I_0(q)}. \]
    由除子方程，左边也等于 \(J(qe^{I_1/I_0}, z)\)，我们也称之为镜像映射。
    因此，我们可以用镜像定理计算 $\star_{\tau}$。\(H^2\) 方向上的量子联络化为常微分方程
    \[ (H+zD) S_{\tau}(z)^* = S_{\tau}(z)^* \cdot I_{11} H \star_{\tau}, \]
    其中 $D \coloneqq q \odv{}{q}$（这里坐标为 $\log q$），$I_{11} \coloneqq 1 + D \ab(\frac{I_1(q)}{I_0(q)})$。利用镜像定理及 Zagier--Zinger~\cite{ZagierZinger2008} 的结果进行显式计算可得
    \begin{align*}
        I_{11} H + \cdots &= S_{\tau}(z)^* (I_{11} H \star_{\tau} \1) \\
        \ab( 1 + D\ab(\frac{\frac{I_1}{I_0} + D \ab(\frac{I_2}{I_0})}{I_{11}}) ) H^2 + \cdots &= S_{\tau}(z)^* (I_{11} H \star_{\tau} H) \\
        I_{11} H^3 &= S_{\tau}(z)^* (I_{11}H \star_{\tau} H^2).
    \end{align*}
    由于 $S_{\tau}$ 以恒同映射起首，这就计算出了 $H \star_{\tau}$。
\end{exm}

\subsubsection{振荡积分}%
\label{ssub:Oscillating integrals}

现在我们重新解释亏格零镜像定理；在讨论高亏格不变量时将回到这一解释。特别地，我们将把 I-函数解释为 \(X\) 的某个镜像族上的积分，对我们而言该镜像族是环面簇。设
\[ 0 \to \Lambda \to \Z^r \to N \to 0 \]
为 \(X\) 的扇序列。施加 \(\Hom(-, \C^{\times})\)，得到正合列
\[ 1 \to (\C^{\times})^{\dim X} \to (\C^{\times})^r \to \Hom(\Lambda, \C^{\times}) \to 1. \]
设映射 \(\Lambda \to \Z^r\) 由矩阵 \((m_{ij})_{i=1,\ldots,r}^{j = 1, \ldots, \rho(X)}\) 给出。考虑与 \(X\) 镜像的簇
\[ E_q = \ab(\prod_{i=1}^r x_j^{m_{ij}} = q_i \middle\vert i=1, \ldots, \rho(X)) \]
并考虑超势 \(f = x_1 + \cdots + x_r\)。若希望考虑等变理论，其中等变参数 \(\lambda_1, \ldots, \lambda_r\) 作用于 \(X\) 的各个齐次坐标，则考虑 \(f_{\lambda} = \sum_{i=1}^r (x_i - \lambda_i \log x_i)\)。

\begin{exm}
    当 \(X = \P^n\) 时，\(X\) 的镜像是超势 \(\sum_{i=1}^{n+1} (x_i - \lambda_i \log x_i)\)，定义在簇 \(E_q = (x_1 \cdots x_{n+1} = q)\) 上。
\end{exm}

现在我们选取合适的 Lefschetz 闭链 \(\Gamma_{\alpha}\)，位于（等变）超势的每个临界点 \((x_{\alpha})\) 附近，并考虑积分
\[ I_{\alpha} = \int_{\Gamma_{\alpha}} e^{f_{\lambda}/z} \prod_{i=1}^r \frac{\d{x_i}}{x_i}. \]
在 I-函数方面，这相当于将表达式中的每个 \(D_i\) 替换为 \(D_i - \lambda_i\)（即打开等变参数）；在 J-函数方面，则相当于考虑 \(X\) 的完全等变 GW 理论。在亏格零时，\(\lambda_i \to 0\) 的极限存在，因此我们可以放心地由下述结果恢复原来的镜像定理：

\begin{thm}[{\cite{Givental1996}}]
    该积分被与 I-函数相同的 Picard–Fuchs 方程所零化，且除各种前因子外，它沿负实轴 \(z \to 0^-\) 时的渐近展开与 I-函数在 \(X\) 上对应于临界点 \(( x_{\alpha} )\) 的 T-不动点处的限制一致。
\end{thm}


\subsubsection{拉格朗日锥}%
\label{ssub:The Lagrangian cone}

\(X\) 的量子上同调是\textit{弗罗贝尼乌斯流形}的一个例子；稍后我们将看到，任何上同调场论都能给出这样的流形。记 \(V = H^*(X)\)，设 \((e_{\mu})\) 为 \(V\) 的一组基，于是可写 $\bv = \sum_{\mu,n} t^{\mu}_n e_{\mu} z^n \in V\ps{z}$。
弗罗贝尼乌斯流形的结构来自函数
\[
    \mc{F}_0(\bv) = \sum_{n, \beta, \mu} \frac{q^{\beta}}{n!} \ab<\bv(\psi), \ldots, \bv(\psi)>_{0,n,\beta},
\]
它也称为 \(X\) 的\textit{亏格零后代势}。它满足弦方程、膨胀子方程以及拓扑递归关系：
\begin{align*}
    \pdv{}{t_0^1} \mc{F}_0 &= \frac{1}{2} (\bv_0, \bv_0) + \sum_{n=0}^{\infty} \sum_{\mu} t_{n+1}^{\mu} \pdv{}{t_n^{\mu}} \mc{F}_0; \\
    \pdv{}{t_1^1} \mc{F}_0 &= \sum_{n=0}^{\infty} \sum_{\mu} t_n^{\mu} \pdv{}{t_n^{\mu}}\mc{F}_0 - 2 \mc{F}_0; \\
    \pdv{}{t_{k+1}^{\alpha},t_m^{\beta},t_n^{\gamma}} \mc{F}_0 &= \sum_{\mu,\nu} \pdv{}{t_k^{\alpha},t_0^{\mu}} \mc{F}_0 \cdot \eta^{\mu\nu} \cdot \pdv{}{t_0^{\nu},t_{m}^{\beta},t_n^{\gamma}} \mc{F}_0.
\end{align*}
为叙述下述结果，考虑带辛形式
\[ \omega(\mbf{f}(z), \mbf{g}(z)) \coloneqq \on{Res}_{z=0}(\mbf{f}(-z), \mbf{g}(z)). \]
的无穷维向量空间 $V\ls{z^{-1}}$（或其某个完备化）。我们还考虑新变量\footnote{这称为\textit{膨胀子平移}。} $\bq = \bv - z$，于是 $\mc{F}_0$ 现在是 $\bq = -z$ 附近的函数。

\begin{thm}[Givental {\cite{Givental2004}}]
    $\mc{F}_0$ 满足上述偏微分方程组，当且仅当 $\mc{L}$（即 $\d{\mc{F}_0}$ 的图像）是顶点位于 $\bq = 0$ 的拉格朗日锥，且其切空间 $L$ 与 $\mc{L}$ 恰好沿 $zL \subset L$ 相切。
\end{thm}

我们可以通过以下步骤恢复拉格朗日锥 $\mc{L}$。先求量子联络的一个基本解
\[ S = \mr{Id} + S_1 z^{-1} + S_2 z^{-2} + \cdots \]
它满足方程
\[ z \pdv{}{t^{\mu}}S = e_{\mu} \star S. \]
然后令 $J \coloneqq z S^*(z) \1$，即可按以下步骤恢复拉格朗日锥：
\begin{itemize}
    \item 导数 $\pdv{}{t^{\mu}} J$ 构成 $L \cap V\ps{z^{-1}}$ 的一组基；
    \item 这意味着 $z \pdv{}{t^{\mu},t^{\nu}} J \in L \cap V \ps{z^{-1}}$。将它们用一阶导数表出，并利用 $J$ 求解量子联络这一事实，我们恢复弗罗贝尼乌斯结构，从而恢复拉格朗日锥。
\end{itemize}

由此，我们可以将亏格零镜像定理解释为：
\begin{thm}[{\cite{Givental1998}}]
    I-函数 \(I_{\mbf{t}}(q,z)\) 位于 \(X\) 的拉格朗日锥上。
\end{thm}




\subsection{曲线模空间与上同调场论}%
\label{sub:Moduli of curves and CohFTs}

上同调场论的概念由 Kontsevich--Manin~\cite{KM1994} 引入。

我们用 $\Mbar_{g,n}$ 记亏格 $g$、带 $n$ 个标记点的稳定曲线的模空间。它非空当且仅当 $2g-2+n > 0$，此时它是维数为 $3g-3+n$ 的 Deligne–Mumford 叠。族 $\Mbar_{g,n}$ 具有组合结构，它由一族态射给出。
\begin{itemize}
    \item 存在粘合态射
        \[q \colon \Mbar_{g,n+1} \times \Mbar_{h, m+1} \to \Mbar_{g+h,n+m}, \]
        它取两条曲线并沿最后一个标记点将它们粘合以形成一个结点；
    \item 存在自粘合态射 \[s \colon \Mbar_{g-1,n+2} \to \Mbar_{g,n},\]
        它粘合最后两个标记点。
\end{itemize}
当然还有另一个有趣的映射，即遗忘映射
\[ \pr \colon \Mbar_{g,n+1} \to \Mbar_{g,n} \]
它删除最后一个标记点然后再作稳定化。现在我们可以定义我们感兴趣的枚举理论，其中包括格罗莫夫–威滕理论。

\begin{defn}
    给定向量空间 $V$，其上带非退化对称双线性型 $\eta$ 及单位元 $\1 \in V$，一个\textit{上同调场论（CohFT）} $\Omega$ 是指 $V$ 上的一族映射 $\Omega_{g,n}$，其中每个都是 $S_n$-等变线性映射
    \[ \Omega_{g,n} \colon V^{\otimes n} \to H^*(\Mbar_{g,n}) \]
    满足基本恒等式
    \[ \Omega_{0,3}(\1, u,v) = \eta(u,v) \]
    及以下组合恒等式（取 $e_{\mu}$ 为 $V$ 的基，使得 $e_1 = \1$，$e^{\mu}$ 为对偶基）：
    \begin{align*}
        q^* \Omega_{g+h,n+m}(\bv_1, \bv_2) &= \sum_{\mu} \Omega_{g,n+1}(\bv_1, e_{\mu}) \cdot \Omega_{h,m+1}(  \bv_2, e^{\mu} ); \\
        s^* \Omega_{g,n}(\bv) &= \sum_{\mu} \Omega_{g,n+2}(\bv, e_{\mu}, e^{\mu}).
    \end{align*}
    若此外还满足恒等式
    \[ \pr^*\Omega_{g,n}(\bv) = \Omega_{g,n+1}(\bv,\1) \]
    则称该 CohFT 具有\textit{平坦单位}（或满足\textit{弦方程}）。
\end{defn}
以上均可推广到超向量空间，但为简单起见我们不处理这种情形。

\begin{exm}
    CohFT 最重要的例子是光滑射影簇 $X$ 的格罗莫夫–威滕理论。这里回顾一下，稳定映射的源是一条 \textbf{预稳定}曲线，因此存在稳定化态射
    \[ \st \colon \Mbar_{g,n}(X,\beta) \to \Mbar_{g,n} \]
    它遗忘映射并对曲线作稳定化。于是在 Novikov 环上工作，取 $V = H^*(X)$，$\eta$ 为庞加莱配对，$\1$ 为 $X$ 的基本类。线性映射 $\Omega^X_{g,n}$ 由下式给出
    \[ \Omega_{g,n}^X(\btau) \coloneqq \sum_{\beta} q^{\beta} \st_* \ab(\prod_{i=1}^n \ev_i^*(\tau_i) \cap [\Mbar_{g,n}(X,\beta)]^{\vir}), \]
    其中 $\ev_i \colon \Mbar_{g,n}(X,\beta) \to X$ 是第 $i$ 个求值映射。
\end{exm}

\begin{exm}
    设 $\pi \colon \mc{C}_{g,n} \to \Mbar_{g,n}$ 为万有曲线。则\textit{霍奇丛}是向量丛
    \[ \E_{g,n} \coloneqq \pi_* \omega_{\pi}. \]
    再考虑带通常配对的向量空间 $V = \C$，并按下式定义 CohFT $\Omega^{\E}$：
    \[ \Omega^{\E}_{g,n} \coloneqq c(\E_{g,n}). \]
    粘合公理由等式 $q^* \E_{g+h} = \E_g \oplus \E_h$ 与正合列
    \[ 0 \to \E_{g-1} \to \E_g \to \mc{O} \to 0. \]
    满足。这里由于亏格零分支对 $\omega$ 的整体截面没有贡献，且 $\E$ 不依赖于标记点，该 CohFT 具有平坦单位。
\end{exm}

\begin{rmk}
    稍后我们将看到，由量子黎曼–罗赫定理，这与一点的 GW CohFT 相关。
\end{rmk}

给定 CohFT $\Omega$，将类 $\Omega_{g,n}(\bv)$ 与 $\Mbar_{g,n}$ 上的各种上同调配对即可得到不变量。其中最值得考虑的是类
\[ \bar{\psi}_i \coloneqq c_1( s_i^* \omega_{\pi} ), \]
这里 $s_i \colon \Mbar_{g,n} \to \mc{C}_{g,n}$ 是第 $i$ 个标记点。

\begin{rmk}
    在格罗莫夫–威滕理论中，这些称为\textit{先祖}。还有后代，它由同样的公式给出，只是用预稳定曲线的模空间 $\mf{M}_{g,n}$ 代替 $\Mbar_{g,n}$，其中我们将上面定义的态射 $\st$ 分解为
    \[ \Mbar_{g,n}(X,\beta) \xrightarrow{\text{forget map}} \mf{M}_{g,n} \xrightarrow{\text{stabilize}} \Mbar_{g,n}. \]
    后代类记作 $\psi_i$。
\end{rmk}

\begin{exm}
    我们将计算一个将在卡拉比–丘三维簇的 GW 理论中出现的不变量。对任意 CohFT $\Omega$，考虑不变量
    \[ \ab<\1\bar{\psi}_1>_{1,1}^{\Omega}\coloneqq \int_{\Mbar_{1,1}} \Omega_{1,1}(\1)\bar{\psi}_1. \]
    利用第二个粘合方程，$\Omega_{1,1}(\1)$ 的零次部分等于
    \[ \sum_{\mu} \Omega_{0,3}(\1, e_{\mu}, e^{\mu}) = \sum_{\mu} \eta(e_{\mu}, e^{\mu}), \]
    这给出（分次）维数 $\chi(V)$，也就是 $V$ 的（分次）维数（例如，若 \(V = H^*(X)\)，其中 \(X\) 为卡拉比–丘三维簇，则这是 \(X\) 的拓扑欧拉示性数）。于是我们得到
    \[ \ab<\1\bar{\psi}_1>_{1,1}^{\Omega} = \frac{\dim V}{24}. \]
\end{exm}

\begin{exm}
    考虑一点的 GW CohFT。它由带通常配对的 $V = \C$ 及公式
    \[ \Omega_{g,n} = 1. \]
    给出。于是我们定义不变量
    \[ \ab<\bar{\psi}_1^{a_1} \cdots \bar{\psi}_n^{a_n}>_{g,n} \coloneqq \int_{\Mbar_{g,n}} \Omega_{g,n} \cdot \bar{\psi}_1^{a_1} \cdots \bar{\psi}_n^{a_n}. \]
    \begin{thm}[Kontsevich {\cite{Kontsevich1992}}]
        函数
        \[ \mc{D} \coloneqq \exp \ab(\sum_{g,n} \frac{\hbar^{g-1}}{n!}\sum_{a_1 + \cdots + a_n = 3g-3+n} \ab<\bar{\psi}_1^{a_1} \cdots \bar{\psi}_n^{a_n}>_{g,n} t_{a_1}\cdots t_{a_n}) \]
        被以下算子零化：
        \begin{align*}
            L_{-1} \coloneqq{}& -\pdv{}{t_0} + \frac{\hbar^{-1}}{2} t_0^2 + \sum_{i=0}^{\infty} t_{i+1} \pdv{}{t_i};\\
            L_0 \coloneqq{}& -\frac{3}{2} \pdv{}{t_1} + \sum_{i=0}^{\infty} \frac{2i+1}{2} t_i \pdv{}{t_i} + \frac{1}{16};\\
            L_n \coloneqq{}& -\frac{(2n+3)!!}{2^{n+1}} \pdv{}{t_{n+1}} + \sum_{i=0}^{\infty} \frac{(2i+2n+1)!!}{(2i-1)!! 2^{n+1}} t_i \pdv{}{t_{i+n}} \\
            &+ \frac{\hbar}{2} \sum_{i=0}^{n-1} \frac{(2i+1) (2i-1) \cdots (2i+1-2n)}{2^{n+1}} \pdv{}{t_i,t_{n-1-i}}.
        \end{align*}
    \end{thm}
\end{exm}

\begin{rmk}
    这一结果（一点的 Virasoro 约束）等价于 $\mc{D}$ 是 KdV 族的 \(\tau\)-函数，并已被多位作者推广。
\end{rmk}



